\documentclass[11pt,letterpaper]{article}

\usepackage[margin=1in]{geometry}
\usepackage[T1]{fontenc}
\usepackage{lmodern}
\usepackage{microtype}
\usepackage{parskip}
\usepackage{enumitem}
\usepackage{booktabs}
\usepackage{titlesec}
\usepackage[hidelinks,colorlinks=true,linkcolor=black,urlcolor=black,citecolor=black]{hyperref}
\usepackage{xcolor}
\usepackage{fancyhdr}

\titleformat{\section}{\Large\bfseries}{\thesection}{1em}{}
\titleformat{\subsection}{\large\bfseries}{\thesubsection}{1em}{}

\pagestyle{fancy}
\fancyhf{}
\fancyhead[L]{\small Theseus Auto-Paper --- bayesian-update-12a26716}
\fancyhead[R]{\small\thepage}
\renewcommand{\headrulewidth}{0.4pt}

\setlength{\parindent}{0pt}

\newcommand{\todomark}[1]{\textcolor{red}{\textbf{[TODO: #1]}}}
\newcommand{\rowref}[1]{\texttt{\small[row:#1]}}

\title{\textbf{Calibrated narrowing under uncertainty: a firm cluster}\\[6pt]\large Cluster bayesian-update-12a26716}
\date{May 2026}
\author{Theseus Auto-Paper Generator}

\begin{document}

\maketitle
\thispagestyle{empty}

\begin{center}
\fbox{%
  \begin{minipage}{0.92\textwidth}
    \centering\small
    \textbf{Disclosure:} machine-drafted, founder-reviewed.\\
    Lead conclusion: \texttt{conclusion-calibration-over-coverage-1}.\\
    Methodology root: \texttt{profile-calibration-over-coverage-1}
    (bayesian-update).
  \end{minipage}%
}
\end{center}

\vspace{1.5em}

\section*{Abstract}
This paper distills 3 firm conclusion(s) anchored on conclusion-calibration-over-coverage-1, all sharing the methodology root "Calibrated narrowing under uncertainty". The cluster's claims are stress-tested against 2 resolved forecast(s) drawn from the firm's calibration ledger. Lead claim: A calibrated narrow claim beats a confident broad claim when the firm must commit; we will revisit if a resolved forecast contradicts the narrowing.

\section{Introduction}
The Theseus engine routinely produces firm conclusions, attaches them to a methodology profile, and tracks how those conclusions perform when their predictive consequences are settled in public forecast markets. This paper presents one such cluster of 3 conclusion(s), drawn together by their shared methodology root and by at least one resolved forecast that touches the cluster. The cluster's lead claim states: "A calibrated narrow claim beats a confident broad claim when the firm must commit; we will revisit if a resolved forecast contradicts the narrowing.".

\subsection*{Conclusions in this cluster}
\begin{itemize}[leftmargin=1.5em]
  \item \textbf{conclusion-calibration-over-coverage-1} \rowref{conclusion:conclusion-calibration-over-coverage-1} --- A calibrated narrow claim beats a confident broad claim when the firm must commit; we will revisit if a resolved forecast contradicts the narrowing.
  \item \textbf{conclusion-calibration-over-coverage-2} \rowref{conclusion:conclusion-calibration-over-coverage-2} --- Reliability priors should decay toward the base rate when the regime shifts, by 2026 Q4 at the latest.
  \item \textbf{conclusion-calibration-over-coverage-3} \rowref{conclusion:conclusion-calibration-over-coverage-3} --- When forced to choose, the firm publishes the calibrated claim and files the broad one as an open question.
\end{itemize}

\section{Methods}
The cluster shares a single methodology root,
\textit{Calibrated narrowing under uncertainty} (\texttt{profile-calibration-over-coverage-1}, pattern \emph{bayesian-update}).

\textbf{Summary.} Resolve a broad question into the narrowest claim the evidence actually licenses, then widen only as resolved forecasts justify it.

\subsection*{Reasoning moves}
\begin{itemize}[leftmargin=1.5em]
  \item anchor on the base rate before reading the signal
  \item prefer the narrow calibrated claim to the broad confident one
  \item widen the claim only when a resolved forecast licenses it
\end{itemize}
\subsection*{Assumptions}
\begin{itemize}[leftmargin=1.5em]
  \item signals are conditionally independent given the regime
  \item the base rate is stationary over the forecast horizon
\end{itemize}
\section{Results}
The cluster's conclusions were stress-tested against
2 resolved forecast(s) from the firm's
calibration ledger. Each row below resolves to a database record.

\begin{center}
\begin{tabular}{@{}p{0.30\textwidth} p{0.10\textwidth} p{0.12\textwidth} p{0.12\textwidth} p{0.20\textwidth}@{}}
\toprule
\textbf{Headline} & \textbf{Outcome} & \textbf{Stated} & \textbf{Brier} & \textbf{Row} \\
\midrule
Calibrated narrowing under uncertainty: claim 1 holds & YES & 0.70 & 0.090 & \rowref{forecast\_resolution:prediction-calibration-over-coverage-1} \\
Calibrated narrowing under uncertainty: claim 2 holds & YES & 0.70 & 0.090 & \rowref{forecast\_resolution:prediction-calibration-over-coverage-2} \\
\bottomrule
\end{tabular}
\end{center}

Across the 2 fully-scored forecast(s) above, the firm achieved a mean Brier score of 0.090. Each contributing row is auditable via its \textbackslash{}rowref\{...\} marker.
\section{Discussion}
Reviewer firm-peer-review on conclusion-calibration-over-coverage-1: verdict revise (confidence 0.78).
\subsection*{Peer-review findings}
\begin{itemize}[leftmargin=1.5em]
  \item \textbf{minor} --- One supporting signal rests on a single source; triangulate before widening the claim.
        \rowref{review\_report:review-calibration-over-coverage-seed}
\end{itemize}
\section{Limits and failure modes}
The methodology root carries a registered failure-mode catalog. Every
limit below is an entry from the firm's catalog, not a synthesized
observation.

\begin{itemize}[leftmargin=1.5em]
  \item regime change invalidates the stationarity assumption
  \item common-mode bias across correlated signals
  \item narrowing past the point the evidence supports
\end{itemize}
\section{References}
\begin{enumerate}[leftmargin=2em]
  \item Firm source memo: Calibrated narrowing under uncertainty, May 2026 \rowref{artifact:artifact-calibration-over-coverage-source-memo}
\end{enumerate}

\vfill
\begin{center}\small
\textit{This paper was machine-drafted by Theseus from the firm's existing
artifacts. Every numerical claim resolves to a row in the firm's
database (see \rowref{...} markers). The .tex source is the authoritative
artifact; the PDF is a build product.}
\end{center}

\end{document}
